\subsection{Pohyb a konstantní rychlost}

Koule se pohybuje konstantní rychlostí na ose X pomocí \texttt{Rigidbody2D}. V každém snímku se nastavuje \texttt{velocity.x = moveSpeed}, což zajišťuje konstantní pohyb doprava bez ohledu na kolize nebo jiné faktory. Tento přímý přístup k rychlosti je jednodušší a spolehlivější než použití sil.

Rigidbody je nastaven na \texttt{freezeRotation}, takže se koule neotáčí vlivem fyziky. Rotace animace koule se řídí samostatně v závislosti na rychlosti pohybu. Toto oddělení fyzikální rotace od vizuální animace umožňuje přesnější kontrolu nad prezentací.

\begin{lstlisting}[language=csharp,caption={Listing 14: Pohyb bowlingové koule}]
void Awake()
{
    rb = GetComponent<Rigidbody2D>();
    rb.freezeRotation = true;
}

void FixedUpdate()
{
    if (!_canMove) return;
    rb.linearVelocity = new Vector2(moveSpeed, rb.linearVelocity.y);
    if (animator != null)
        animator.speed = Mathf.Max(0.1f, rb.linearVelocity.magnitude / 5f);
}
\end{lstlisting}
